\documentclass[a4paper]{article} %

\date{} %

% Please, do not change any of the above lines

\usepackage{amssymb} % add other LaTeX packages you need here.
\usepackage{doi}

\begin{document}
\setcounter{page}{1} % Please, do not delete this line

% Your title
\title{Exploration of non-classical negation axioms}

% Authors
\author{Mike Behrisch\\
        TU Wien\\ \\ 
       \texttt{behrisch@logic.at}
       }%

\maketitle

% The abstract
We explore the implicational interconnections between different types
of lattice based non-classical logic algebras. That is, we consider
structures of the form
$\mathbf{L} = \langle L; \wedge, \vee, 0,1, \overline{\rule{0pt}{1ex}~}\rangle$ with
the background assumption that $\langle L; \wedge,\vee,0,1\rangle$ is a
bounded lattice.
Our investigation focusses on the additional
axioms the `negation' (or `complementation') operator
$x\mapsto\overline{x}$ may fulfil, and on the implications between those
properties. In this respect we consider several generalisations of
orthocomplemented lattices (as used in quantum logic) and of
pseudocomplemented lattices (as appearing, \emph{e.g.}, in Heyting
algebras when dealing with intuitionistic logic).
The properties of negation we study are mainly taken from the ones
surveyed in~\cite{HollidayAFundamentalNonClassicalLogic2023}, without
any claims to represent the available literature with any degree of
completeness. A different set of negation axioms was examined
in~\cite{DunnZhouNegationInTheContextOfGaggleTheory}, but the final
picture we get of the situation is richer than both of these works and
more along the lines
of~\cite{DauImplicationsComplementationPropertiesFiniteLattices}, where
again other properties were looked at.

\begin{thebibliography}{9}
\bibitem{DauImplicationsComplementationPropertiesFiniteLattices}
F.~Dau.
\newblock \emph{Implications of properties concerning complementation in finite lattices,}
\newblock In Contributions to General Algebra \textbf{12}, ({V}ienna, 1999), 145--154.
          Heyn, Klagenfurt, 2000.

\bibitem{DunnZhouNegationInTheContextOfGaggleTheory}
J.\,M.~Dunn and Ch.~Zhou.
\newblock \emph{Negation in the context of gaggle theory,}
\newblock Studia Logica \textbf{80} 235--264 (2005)
\newblock \doi{10.1007/s11225-005-8470-y}

\bibitem{HollidayAFundamentalNonClassicalLogic2023}
W.\,H.~Holliday.
\newblock \emph{A fundamental non-classical logic,}
\newblock Logics \textbf{1} 36--79 (2023)
\newblock \doi{10.3390/logics1010004}
\end{thebibliography}

\end{document}
